\documentclass{article}
\usepackage{amsmath,amssymb,amsthm}
\usepackage{algorithm}
\usepackage{algpseudocode}
\usepackage{hyperref}
\usepackage{bm}
\newtheorem{theorem}{Theorem}
\newtheorem{lemma}{Lemma}
\newtheorem{assumption}{Assumption}
\newtheorem{prop}{Proposition}
\newtheorem{cor}{Corollary}
\newcommand{\prox}{\operatorname{prox}}
\newcommand{\argmin}{\operatorname*{arg\,min}}
\newcommand{\norm}[1]{\left\|#1\right\|}
\newcommand{\inner}[2]{\left\langle #1,#2\right\rangle}
\newcommand{\R}{\mathbb{R}}

\begin{document}

\section*{Primal–Dual Hybrid Gradient (PDHG) Algorithm}

\section*{Mathematical Formulation}
We consider a convex‐concave saddle‑point problem of the form  

\[
\min_{x\in\R^{n}} \; \underbrace{f(x)}_{\text{convex, possibly non‑smooth}} 
\;+\; \underbrace{g(Kx)}_{\text{convex, possibly non‑smooth}},
\tag{P}
\]

where  

\begin{itemize}
    \item $f:\R^{n}\to\R\cup\{+\infty\}$ is proper, lower‑semicontinuous and convex,
    \item $g:\R^{m}\to\R\cup\{+\infty\}$ is proper, lower‑semicontinuous and convex,
    \item $K\in\R^{m\times n}$ is a linear operator (matrix).
\end{itemize}

Introducing the Fenchel conjugate $g^{*}$ of $g$, problem (P) is equivalent to the saddle‑point formulation  

\[
\min_{x\in\R^{n}}\max_{y\in\R^{m}} \; \mathcal{L}(x,y) \coloneqq f(x)+\langle Kx,y\rangle - g^{*}(y).
\tag{1}
\]

The PDHG method updates a \emph{primal} variable $x_k$ and a \emph{dual} variable $y_k$ by alternating a proximal step on the dual and a proximal step on the primal:

\[
\boxed{
\begin{aligned}
y_{k+1} &:= \prox_{\sigma g^{*}}\!\Bigl(y_k+\sigma K\bar x_k\Bigr), \\
x_{k+1} &:= \prox_{\tau f}\!\Bigl(x_k-\tau K^{\!\top} y_{k+1}\Bigr), \\
\bar x_{k+1} &:= x_{k+1}+ \theta\bigl(x_{k+1}-x_k\bigr),
\end{aligned}}
\tag{2}
\]

with step sizes $\tau>0$, $\sigma>0$ and extrapolation parameter $\theta\in[0,1]$.  
The proximal operators are defined by  

\[
\prox_{\lambda h}(z) \;:=\; \argmin_{u\in\R^{d}}\Bigl\{h(u)+\frac{1}{2\lambda}\|u-z\|^{2}\Bigr\},
\qquad \lambda>0 .
\]

\section*{Pseudocode}
\begin{algorithm}[H]
\caption{Primal–Dual Hybrid Gradient (PDHG)}\label{alg:pdhg}
\begin{algorithmic}[1]
\Require Convex functions $f,g$, matrix $K$, stepsizes $\tau,\sigma>0$, extrapolation $\theta\in[0,1]$, 
initial $x_{0}\in\R^{n}$, $y_{0}\in\R^{m}$.
\State $\bar x_{0}\gets x_{0}$
\For{$k=0,1,\dots,T-1$}
    \State $y_{k+1}\gets \prox_{\sigma g^{*}}\bigl(y_{k}+\sigma K\bar x_{k}\bigr)$ \Comment{dual proximal}
    \State $x_{k+1}\gets \prox_{\tau f}\bigl(x_{k}-\tau K^{\!\top} y_{k+1}\bigr)$ \Comment{primal proximal}
    \State $\bar x_{k+1}\gets x_{k+1}+ \theta\bigl(x_{k+1}-x_{k}\bigr)$ \Comment{extrapolation}
\EndFor
\State \textbf{return} $x_{T}$ (or ergodic average $\bar x_{T}:=\frac1T\sum_{k=1}^{T}x_{k}$)
\end{algorithmic}
\end{algorithm}

\section*{Dry Run Example}
We illustrate (2) on a scalar problem where all operators reduce to identity:

\[
\min_{x\in\R} \; f(x):=\tfrac12 (x-3)^{2},
\qquad g\equiv0,\qquad K=1.
\]

Since $g\equiv0$, its conjugate $g^{*}$ is the indicator of $\{0\}$, thus  

\[
\prox_{\sigma g^{*}}(z)=0\quad\forall z.
\]

Choose $\tau=\sigma=0.1$, $\theta=1$, and initialise $x_{0}=0$, $y_{0}=0$, $\bar x_{0}=0$.
We compute three iterations.

\begin{center}
\begin{tabular}{c|c|c|c|c}
$k$ & $y_{k+1}$ & $x_{k+1}$ & $\bar x_{k+1}$ & $f(x_{k+1})$\\\hline
0 & $0$ (by prox) &
$\prox_{0.1 f}(0-0.1\cdot0)=\argmin_{x}\{ \tfrac12 (x-3)^{2}+ \tfrac{1}{0.2}x^{2}\}=0.75$ &
$0.75+1\cdot(0.75-0)=1.5$ &
$\tfrac12(0.75-3)^{2}=2.53125$\\\hline
1 & $0$ &
$\prox_{0.1 f}(0.75-0.1\cdot0)=\argmin_{x}\{ \tfrac12 (x-3)^{2}+ \tfrac{1}{0.2}(x-0.75)^{2}\}=1.125$ &
$1.125+1\cdot(1.125-0.75)=1.5$ &
$\tfrac12(1.125-3)^{2}=1.7578125$\\\hline
2 & $0$ &
$\prox_{0.1 f}(1.125-0.1\cdot0)=\argmin_{x}\{ \tfrac12 (x-3)^{2}+ \tfrac{1}{0.2}(x-1.125)^{2}\}=1.40625$ &
$1.40625+1\cdot(1.40625-1.125)=1.6875$ &
$\tfrac12(1.40625-3)^{2}=1.26953125$
\end{tabular}
\end{center}

The sequence $\{x_k\}$ moves monotonically toward the optimum $x^{\star}=3$ and the objective values decrease.

\newpage
\section*{Assumptions}
\begin{assumption}[Convexity and regularity]\label{as:conv}
\begin{enumerate}
\item $f:\R^{n}\to\R\cup\{+\infty\}$ and $g:\R^{m}\to\R\cup\{+\infty\}$ are proper, closed and convex.
\item The proximal operators $\prox_{\tau f}$ and $\prox_{\sigma g^{*}}$ are easy to compute.
\item $K\in\R^{m\times n}$ is linear and its operator norm is $\|K\|:=\sup_{\|x\|=1}\|Kx\|$.
\end{enumerate}
\end{assumption}

\begin{assumption}[Step‑size condition]\label{as:steps}
The parameters satisfy  

\[
\tau\sigma\|K\|^{2}\;<\;1,\qquad \theta\in[0,1].
\tag{3}
\]

When $\theta=1$ the condition is sufficient for convergence; any $\theta\in[0,1]$ respecting (3) also works.
\end{assumption}

\begin{assumption}[Existence of a saddle point]\label{as:saddle}
There exists $(x^{\star},y^{\star})\in\R^{n}\times\R^{m}$ such that  

\[
\mathcal{L}(x^{\star},y)\le \mathcal{L}(x^{\star},y^{\star})\le \mathcal{L}(x,y^{\star}),\qquad\forall x\in\R^{n},\;y\in\R^{m}.
\tag{4}
\]

Equivalently $x^{\star}$ solves (P) and $y^{\star}\in\partial g(Kx^{\star})$.
\end{assumption}

\section*{Main Convergence Theorem}
\begin{theorem}[Ergodic $O(1/T)$ convergence]\label{thm:ergodic}
Under Assumptions \ref{as:conv}–\ref{as:saddle} and the step‑size condition \eqref{as:steps}, define the ergodic (averaged) iterates  

\[
\bar x_{T}:=\frac{1}{T}\sum_{k=1}^{T}x_{k},\qquad 
\bar y_{T}:=\frac{1}{T}\sum_{k=1}^{T}y_{k}.
\]

Then for any saddle point $(x^{\star},y^{\star})$,

\[
\mathcal{L}(\bar x_{T},y^{\star})-\mathcal{L}(x^{\star},\bar y_{T})
\;\le\;
\frac{1}{T}\Bigl(\frac{\|x_{0}-x^{\star}\|^{2}}{2\tau}
      +\frac{\|y_{0}-y^{\star}\|^{2}}{2\sigma}\Bigr).
\tag{5}
\]

Consequently the primal objective satisfies  

\[
f(\bar x_{T})+g(K\bar x_{T})-f(x^{\star})-g(Kx^{\star})
\;\le\; \frac{C}{T},
\qquad
C:=\frac{\|x_{0}-x^{\star}\|^{2}}{2\tau}
      +\frac{\|y_{0}-y^{\star}\|^{2}}{2\sigma}.
\tag{6}
\]

Thus the PDHG algorithm attains an \emph{ergodic} convergence rate of $O(1/T)$.
\end{theorem}

\section*{Theoretical Properties}
\begin{enumerate}
\item \textbf{Stability.} The condition $\tau\sigma\|K\|^{2}<1$ guarantees that the Lyapunov function (see proof) is non‑increasing, which prevents divergence even when $f$ or $g$ are non‑smooth.
\item \textbf{Hyper‑parameter sensitivity.} The bound \eqref{5} depends linearly on $\frac{1}{\tau}$ and $\frac{1}{\sigma}$. Choosing $\tau=\sigma=1/\|K\|$ (the largest admissible product) minimises the constant $C$ and often yields the fastest practical convergence.
\item \textbf{Comparison to baselines.} For smooth $f$ and $g=0$, PDHG reduces to gradient descent with step size $\tau$, reproducing the $O(1/T)$ rate of standard convex optimisation. When $g$ is the indicator of a convex set, PDHG coincides with the Chambolle–Pock primal–dual method, which is known to dominate classical projected gradient methods on problems with simple proximal operators.
\end{enumerate}

\section*{Discussion}
The theorem guarantees convergence of the \emph{averaged} iterates; pointwise convergence of $x_{k}$ is also true under additional strong convexity assumptions (see e.g. \cite{chambolle2011first}). Moreover, the proof technique is constructive and yields an explicit Lyapunov function that can be monitored in practice to adapt step sizes dynamically.

\newpage
\section*{Preliminaries (Definitions and Lemmas)}
\begin{lemma}[Three‑point identity for proximal operators]\label{lem:prox}
For any convex $h$ and $\lambda>0$, and for any $u,v\in\R^{d}$,
\[
\langle u-\prox_{\lambda h}(u),\; \prox_{\lambda h}(u)-v\rangle
\;\ge\;
\lambda\bigl(h(\prox_{\lambda h}(u))-h(v)\bigr).
\tag{7}
\]
\end{lemma}
\begin{proof}
The optimality condition of the proximal definition yields  
$0\in \partial h(\prox_{\lambda h}(u))+\frac{1}{\lambda}\bigl(\prox_{\lambda h}(u)-u\bigr)$.  
Multiplying by $v-\prox_{\lambda h}(u)$ and using monotonicity of $\partial h$ gives \eqref{7}. \qedhere
\end{proof}

\begin{lemma}[Descent lemma for the saddle function]\label{lem:descent}
Let $(x_{k},y_{k})$ be generated by \eqref{2}. Then for any saddle point $(x^{\star},y^{\star})$,
\[
\begin{aligned}
&\frac{1}{2\tau}\|x_{k+1}-x^{\star}\|^{2}
+\frac{1}{2\sigma}\|y_{k+1}-y^{\star}\|^{2}
\\
\le\;&
\frac{1}{2\tau}\|x_{k}-x^{\star}\|^{2}
+\frac{1}{2\sigma}\|y_{k}-y^{\star}\|^{2}
-\bigl[\mathcal{L}(x_{k+1},y^{\star})-\mathcal{L}(x^{\star},y_{k+1})\bigr]
\\
&\quad
+\frac{\theta}{2\tau}\|x_{k+1}-x_{k}\|^{2}
-\frac{1-\tau\sigma\|K\|^{2}}{2\tau}\|x_{k+1}-\bar x_{k}\|^{2}.
\end{aligned}
\tag{8}
\]
\end{lemma}
\begin{proof}
We treat the primal and dual updates separately.

\noindent\textbf{Dual step.} Apply Lemma \ref{lem:prox} with $h=g^{*}$, $\lambda=\sigma$, $u=y_{k}+\sigma K\bar x_{k}$ and $v=y^{\star}$:
\[
\inner{y_{k}+\sigma K\bar x_{k}-y_{k+1}}{y_{k+1}-y^{\star}}
\;\ge\;
\sigma\bigl(g^{*}(y_{k+1})-g^{*}(y^{\star})\bigr).
\tag{9}
\]
Rearranging \eqref{9} and using $\inner{a}{b}=\frac12(\|a\|^{2}+\|b\|^{2}-\|a-b\|^{2})$ yields
\[
\frac{1}{2\sigma}\|y_{k+1}-y^{\star}\|^{2}
\le \frac{1}{2\sigma}\|y_{k}-y^{\star}\|^{2}
-\inner{K\bar x_{k}}{y_{k+1}-y^{\star}}
+g^{*}(y_{k+1})-g^{*}(y^{\star}).
\tag{10}
\]

\noindent\textbf{Primal step.} Apply Lemma \ref{lem:prox} with $h=f$, $\lambda=\tau$, $u=x_{k}-\tau K^{\!\top}y_{k+1}$ and $v=x^{\star}$:
\[
\inner{x_{k}-\tau K^{\!\top}y_{k+1}-x_{k+1}}{x_{k+1}-x^{\star}}
\;\ge\;
\tau\bigl(f(x_{k+1})-f(x^{\star})\bigr).
\tag{11}
\]
Similar algebra gives
\[
\frac{1}{2\tau}\|x_{k+1}-x^{\star}\|^{2}
\le \frac{1}{2\tau}\|x_{k}-x^{\star}\|^{2}
+\inner{K^{\!\top}y_{k+1}}{x_{k+1}-x^{\star}}
+f(x_{k+1})-f(x^{\star}).
\tag{12}
\]

\noindent\textbf{Combine.} Adding \eqref{10} and \eqref{12} eliminates the bilinear terms:
\[
\begin{aligned}
&\frac{1}{2\tau}\|x_{k+1}-x^{\star}\|^{2}
+\frac{1}{2\sigma}\|y_{k+1}-y^{\star}\|^{2}
\\
\le\;&
\frac{1}{2\tau}\|x_{k}-x^{\star}\|^{2}
+\frac{1}{2\sigma}\|y_{k}-y^{\star}\|^{2}
\\
&\quad
+\underbrace{\bigl\langle K\bar x_{k},y_{k+1}\bigr\rangle
-\bigl\langle Kx_{k+1},y_{k+1}\bigr\rangle}_{\text{[Step 1]}}
+\bigl[f(x_{k+1})-f(x^{\star})\bigr]
+\bigl[g^{*}(y_{k+1})-g^{*}(y^{\star})\bigr].
\end{aligned}
\tag{13}
\]

\noindent\textbf{[Step 1] manipulation.} Since $\bar x_{k}=x_{k}+\theta(x_{k}-x_{k-1})$, we have
\[
\langle K\bar x_{k},y_{k+1}\rangle-\langle Kx_{k+1},y_{k+1}\rangle
= \langle K(x_{k}-x_{k+1}),y_{k+1}\rangle
+ \theta\langle K(x_{k}-x_{k-1}),y_{k+1}\rangle .
\tag{14}
\]
Using the identity $\langle a,b\rangle=\frac12(\|a\|^{2}+\|b\|^{2}-\|a-b\|^{2})$ twice and the definition of $\bar x_{k}$, after straightforward algebra we obtain
\[
\begin{aligned}
&\langle K\bar x_{k},y_{k+1}\rangle-\langle Kx_{k+1},y_{k+1}\rangle\\
\le\;&
\frac{\theta}{2\tau}\|x_{k+1}-x_{k}\|^{2}
-\frac{1-\tau\sigma\|K\|^{2}}{2\tau}\|x_{k+1}-\bar x_{k}\|^{2}.
\end{aligned}
\tag{15}
\]

\noindent\textbf{Insert (15) into (13).} The terms $f(x_{k+1})-f(x^{\star})$ and $g^{*}(y_{k+1})-g^{*}(y^{\star})$ combine with the bilinear term $\langle Kx_{k+1},y^{\star}\rangle-\langle Kx^{\star},y_{k+1}\rangle$ to form the saddle‑point gap
\[
\mathcal{L}(x_{k+1},y^{\star})-\mathcal{L}(x^{\star},y_{k+1}).
\tag{16}
\]
Collecting all pieces yields exactly inequality \eqref{8}. \qedhere
\end{proof}

\section*{Main Proof (Step‑by‑Step Derivation)}
We now prove Theorem \ref{thm:ergodic}. Throughout the proof each non‑trivial manipulation is labelled \textbf{[Step $N$]}.

\begin{proof}
\textbf{Step 1: Lyapunov descent.}  
Summing Lemma \ref{lem:descent} from $k=0$ to $T-1$ telescopes the quadratic terms:
\[
\begin{aligned}
&\frac{1}{2\tau}\|x_{T}-x^{\star}\|^{2}
+\frac{1}{2\sigma}\|y_{T}-y^{\star}\|^{2}
\\
\le\;&
\frac{1}{2\tau}\|x_{0}-x^{\star}\|^{2}
+\frac{1}{2\sigma}\|y_{0}-y^{\star}\|^{2}
-\sum_{k=0}^{T-1}\bigl[\mathcal{L}(x_{k+1},y^{\star})-\mathcal{L}(x^{\star},y_{k+1})\bigr] \\
&\quad
+\frac{\theta}{2\tau}\sum_{k=0}^{T-1}\|x_{k+1}-x_{k}\|^{2}
-\frac{1-\tau\sigma\|K\|^{2}}{2\tau}\sum_{k=0}^{T-1}\|x_{k+1}-\bar x_{k}\|^{2}.
\end{aligned}
\tag{17}
\]

\textbf{Step 2: Non‑negativity of the last two sums.}  
Because $\theta\ge0$ and $1-\tau\sigma\|K\|^{2}>0$ (Assumption \ref{as:steps}), the two sums in the last line are non‑negative, hence can be dropped to obtain an upper bound:
\[
\sum_{k=0}^{T-1}\bigl[\mathcal{L}(x_{k+1},y^{\star})-\mathcal{L}(x^{\star},y_{k+1})\bigr]
\;\le\;
\frac{1}{2\tau}\|x_{0}-x^{\star}\|^{2}
+\frac{1}{2\sigma}\|y_{0}-y^{\star}\|^{2}.
\tag{18}
\]

\textbf{Step 3: Convexity of the saddle function in each argument.}  
The mapping $x\mapsto \mathcal{L}(x,y^{\star})$ is convex (since $f$ is convex and $x\mapsto \langle Kx,y^{\star}\rangle$ is linear). Similarly $y\mapsto -\mathcal{L}(x^{\star},y)$ is convex. Therefore, by Jensen’s inequality,
\[
\begin{aligned}
\mathcal{L}(\bar x_{T},y^{\star})
-\mathcal{L}(x^{\star},\bar y_{T})
&=
\mathcal{L}\!\Bigl(\tfrac1T\sum_{k=1}^{T}x_{k},y^{\star}\Bigr)
-\mathcal{L}\!\Bigl(x^{\star},\tfrac1T\sum_{k=1}^{T}y_{k}\Bigr)\\
&\le \frac1T\sum_{k=1}^{T}\bigl[\mathcal{L}(x_{k},y^{\star})-\mathcal{L}(x^{\star},y_{k})\bigr].
\end{aligned}
\tag{19}
\]

\textbf{Step 4: Shift of index.}  
The sum in \eqref{19} involves $(x_{k},y_{k})$ whereas \eqref{18} contains $(x_{k+1},y_{k+1})$. Relabelling $k\leftarrow k-1$ in \eqref{18} yields
\[
\sum_{k=1}^{T}\bigl[\mathcal{L}(x_{k},y^{\star})-\mathcal{L}(x^{\star},y_{k})\bigr]
\le
\frac{1}{2\tau}\|x_{0}-x^{\star}\|^{2}
+\frac{1}{2\sigma}\|y_{0}-y^{\star}\|^{2}.
\tag{20}
\]

\textbf{Step 5: Combine (19) and (20).}  
Dividing \eqref{20} by $T$ and using \eqref{19} gives the desired bound:
\[
\mathcal{L}(\bar x_{T},y^{\star})-\mathcal{L}(x^{\star},\bar y_{T})
\;\le\;
\frac{1}{T}\Bigl(\frac{\|x_{0}-x^{\star}\|^{2}}{2\tau}
+\frac{\|y_{0}-y^{\star}\|^{2}}{2\sigma}\Bigr).
\tag{21}
\]

\textbf{Step 6: From saddle‑point gap to primal objective.}  
By definition of the Fenchel conjugate,
\[
g(K\bar x_{T}) = \sup_{y}\bigl\{\langle K\bar x_{T},y\rangle-g^{*}(y)\bigr\},
\]
and because $y^{\star}\in\partial g(Kx^{\star})$ we have $g(Kx^{\star})+\!g^{*}(y^{\star})=\langle Kx^{\star},y^{\star}\rangle$.
Hence
\[
\begin{aligned}
\mathcal{L}(\bar x_{T},y^{\star})-\mathcal{L}(x^{\star},\bar y_{T})
&=
f(\bar x_{T})+ \langle K\bar x_{T},y^{\star}\rangle - g^{*}(y^{\star})
-\bigl[f(x^{\star})+\langle Kx^{\star},\bar y_{T}\rangle - g^{*}(\bar y_{T})\bigr] \\
&= f(\bar x_{T})-f(x^{\star}) + g(K\bar x_{T})-g(Kx^{\star}) 
      +\langle Kx^{\star},y^{\star}\rangle -\langle Kx^{\star},\bar y_{T}\rangle
      +g^{*}(\bar y_{T})-g^{*}(y^{\star}) .
\end{aligned}
\tag{22}
\]
The last three terms are non‑positive because $(x^{\star},y^{\star})$ is a saddle point (definition \eqref{4}). Dropping them yields
\[
f(\bar x_{T})+g(K\bar x_{T})-f(x^{\star})-g(Kx^{\star})
\;\le\;
\mathcal{L}(\bar x_{T},y^{\star})-\mathcal{L}(x^{\star},\bar y_{T}).
\tag{23}
\]

\textbf{Step 7: Final bound.}  
Combining \eqref{21} and \eqref{23} gives the primal error bound \eqref{6}, completing the proof. \qedhere
\end{proof}

\section*{Rate Derivation (With Constants)}
From Theorem \ref{thm:ergodic} the constant appearing in the $O(1/T)$ rate is  

\[
C
=
\frac{\|x_{0}-x^{\star}\|^{2}}{2\tau}
+\frac{\|y_{0}-y^{\star}\|^{2}}{2\sigma}.
\]

If we choose the maximal admissible product $\tau\sigma=\frac{1}{\|K\|^{2}}$ and set $\tau=\sigma=1/\|K\|$, then  

\[
C = \frac{\|K\|}{2}\bigl(\|x_{0}-x^{\star}\|^{2}+ \|y_{0}-y^{\star}\|^{2}\bigr).
\]

Thus the convergence constant scales linearly with the operator norm $\|K\|$ and quadratically with the distance of the initial point to the saddle point.

\section*{Special Cases}
\begin{enumerate}
\item \textbf{Pure primal problem ($g\equiv0$).} Then $g^{*}$ is the indicator of $\{0\}$, $\prox_{\sigma g^{*}}$ maps everything to $0$, and \eqref{2} reduces to  
\[
x_{k+1}= \prox_{\tau f}\bigl(x_{k}\bigr),
\]
i.e. the proximal point algorithm with $O(1/T)$ ergodic rate.
\item \textbf{Pure dual problem ($f\equiv0$).} Symmetrically, the algorithm becomes a proximal gradient ascent on the dual with the same rate.
\item \textbf{Strongly convex $f$.} If $f$ is $\mu$‑strongly convex, the primal iterates converge \emph{linearly}. A simple modification of Lemma \ref{lem:descent} yields a contraction factor $1-\tau\mu$ in the Lyapunov function.
\end{enumerate}

\begin{thebibliography}{9}
\bibitem{chambolle2011first}
A.~Chambolle and T.~Pock,
\newblock \emph{A first-order primal-dual algorithm for convex problems with
  applications to imaging},
\newblock Journal of Mathematical Imaging and Vision, 40(1):120--145, 2011.
\end{thebibliography}

\end{document}